\documentclass[12pt, a4paper]{article}

\title{
    SINOPSIS\\
    MAGANG INDUSTRI SEBAGAI \textit{FRONT-END DEVELOPER}\\
    DI PT INTI UTAMA SOLUSINDO
}
\author{535230139 ALDIAN YOHANES}

% Format Font Calibri
\usepackage{fontspec}
\setmainfont{Calibri}

% Format Page Layout / Margin
\usepackage[
    a4paper, 
    left=4cm, 
    right=3cm, 
    top=4cm, 
    bottom=3cm,
    showframe
]{geometry}


% Default spacing ke 1.5
\usepackage{setspace}
\onehalfspacing 
\setlength{\parindent}{0cm}
\setlength{\parskip}{8pt}

% Import package lainnya (gambar, tabel, daftar isi)
\usepackage[indonesian]{babel}

\begin{document}
\makeatletter
\begin{center}
    \large\textbf{\@title}
\end{center}

\makeatother
PT Inti Utama Solusindo (IUS) merupakan perusahaan yang bergerak di bidang teknologi informasi dan menjadi bagian dari Pharos Group, sebuah grup perusahaan yang beroperasi dalam industri farmasi dan layanan kesehatan di Indonesia. Sejalan dengan transformasi digital pada sektor healthcare, PT IUS berperan sebagai pengembang sistem dan aplikasi digital yang mendukung operasional internal grup maupun layanan eksternal kepada fasilitas kesehatan.\\
Perusahaan ini mengembangkan berbagai platform berbasis web dan mobile, seperti sistem informasi klinik (HealthyOne), platform distribusi farmasi (Hospinet), serta aplikasi pendukung operasional pemasaran dan marketplace internal. Fokus utama kegiatan perusahaan adalah pengembangan dan pemeliharaan sistem informasi terintegrasi yang mendukung efisiensi, akurasi data, dan optimalisasi layanan kesehatan berbasis teknologi.\\
Kegiatan magang industri dilaksanakan di kantor PT Inti Utama Solusindo yang berlokasi di Gedung Century, Jl. Limo No.6/45, RT.14/RW.02, Grogol Utara, Kec. Kebayoran Lama, Kota Jakarta Selatan. Program magang berlangsung selama satu tahun, dimulai pada 12 Januari 2026 hingga 12 Januari 2027.\\
Pelaksanaan kerja dilakukan secara penuh di kantor (Work From Office/WFO) dengan jadwal kerja Senin sampai Jumat pukul 08.00–17.00 WIB. Skema magang ini termasuk dalam program MBKM (Merdeka Belajar–Kampus Merdeka) Fakultas Teknologi Informasi Universitas Tarumanagara dalam semester ke-5 yang dilanjutkan dengan konversi mata kuliah magang pada semester ke-7.\\
Magang industri ini dilaksanakan sebagai bagian dari implementasi program MBKM untuk memberikan pengalaman profesional langsung kepada mahasiswa dalam lingkungan kerja nyata. Tujuan magang secara khusus meliputi peningkatan kompetensi teknis di bidang pengembangan perangkat lunak, terutama Front-End Development menggunakan Flutter dan React, penerapan teori perkuliahan seperti Software Development Life Cycle (SDLC), Object-Oriented Programming (OOP), clean code, dan prinsip SOLID dalam proyek industri nyata, pengembangan kemampuan kolaborasi dalam tim IT yang terdiri dari Project Manager, Quality Assurance, dan developer lainnya, pemahaman sistem manajemen proyek serta version control profesional menggunakan OpenProject dan GitLab, serta peningkatan kesiapan kerja, profesionalisme, dan kemampuan adaptasi terhadap standar industri teknologi informasi.\\
Bidang kegiatan utama PT Inti Utama Solusindo mencakup pengembangan sistem dan aplikasi berbasis teknologi informasi untuk mendukung ekosistem layanan kesehatan dan farmasi. Aktivitas tersebut meliputi pengembangan aplikasi sistem informasi klinik berbasis web dan mobile, pembangunan platform distribusi dan marketplace farmasi digital, pemeliharaan sistem, perbaikan bug, dan pengembangan fitur baru sesuai kebutuhan bisnis, serta integrasi sistem berbasis API dan manajemen database. Sebagai Front-End Developer Intern, fokus pekerjaan adalah pada pengembangan antarmuka pengguna, migrasi aplikasi Flutter ke versi terbaru, serta pengembangan ulang sistem berbasis React dengan pendekatan modular dan reusable component.\\

\end{document}